\begin{figure*}[t]
\centering
\begin{tikzpicture}[
  font=\small,
  stage/.style={draw, rounded corners=2pt, align=center,
                minimum width=20mm, minimum height=8mm},
  skip/.style={draw, rounded corners=2pt, align=center,
               minimum width=20mm, minimum height=8mm, dashed,
               text=gray},
  arrow/.style={-{Latex[length=2mm]}, semithick},
  row/.style={font=\small\bfseries, anchor=east}
]
\node[row] at (-0.5,0) {No PT};
\node[row] at (-0.5,-1.05) {PT, no IT};
\node[row] at (-0.5,-2.10) {Short};
\node[row] at (-0.5,-3.15) {Long bilingual};

\node[stage,fill=stagegray] (base0) at (1,0) {Base};
\node[stage,fill=stagegray] (base1) at (1,-1.05) {Base};
\node[stage,fill=stagegray] (base2) at (1,-2.10) {Base};
\node[stage,fill=stagegray] (base3) at (1,-3.15) {Base};

\node[skip]  (pt0) at (3.6,0) {PT};
\node[stage,fill=stagegray] (it0) at (6.2,0) {IT};
\node[stage,fill=stagegreen] (act0) at (8.8,0) {Action};
\node[stage,fill=stagegreen] (traj0) at (11.4,0) {Trajectory};

\node[stage,fill=longorange] (pt1) at (3.6,-1.05) {PT};
\node[skip]  (it1) at (6.2,-1.05) {IT};
\node[stage,fill=stagegreen] (act1) at (8.8,-1.05) {Action};
\node[stage,fill=stagegreen] (traj1) at (11.4,-1.05) {Trajectory};

\node[stage,fill=shortblue] (pt2) at (3.6,-2.10) {Short PT};
\node[stage,fill=stagegray] (it2) at (6.2,-2.10) {IT};
\node[stage,fill=stagegreen] (act2) at (8.8,-2.10) {Action};
\node[stage,fill=stagegreen] (traj2) at (11.4,-2.10) {Trajectory};

\node[stage,fill=longorange] (pt3) at (3.6,-3.15) {Long PT\\EN+KO};
\node[stage,fill=stagegray] (it3) at (6.2,-3.15) {IT};
\node[stage,fill=stagegreen] (act3) at (8.8,-3.15) {Action};
\node[stage,fill=stagegreen] (traj3) at (11.4,-3.15) {Trajectory};

\foreach \r in {0,1,2,3} {
  \draw[arrow] (base\r) -- (pt\r);
  \draw[arrow] (pt\r) -- (it\r);
  \draw[arrow] (it\r) -- (act\r);
  \draw[arrow] (act\r) -- (traj\r);
}
\end{tikzpicture}
\caption{\textbf{실제로 준비한 네 학습 경로.} 점선 상자는 해당
단계를 건너뜀을 뜻한다. 각 단계의 데이터 구성과 비교 목적은
본문에서 설명한다.}
\label{fig:pipeline}
\end{figure*}
